%!TEX root = ED_Grafo.tex

\section{Rotas}
\begin{frame}
	\frametitle{Caminhos mais Curtos}
	\framesubtitle{Menor Caminho \emph{x} Caminho de menor custo}

	\begin{center}
		\includegraphics[width=.75\textwidth]{Theofilos_monomaxia_Achilleos}\\
		\small Pintura de Theophilos Hatzimihail
	\end{center}
\end{frame}

\begin{frame}[label=notas]
	\begin{description}
		\item[Caminho Euleriano] passa por todas as arestas do grafo \emph{uma} vez.
		\item[Caminho Hamiltoniano] passa por todos os vértices \emph{uma} vez.
	\end{description}
	\vfill\pause
	\begin{description}
		\item[Ciclo Euleriano] Circuito! Como verificar se existe?
		\item[Ciclo Hamiltoniano] Existe um algoritmo eficiente para encontrar um?
	\end{description}
\end{frame}

\subsection{As Pontes de Königsberg}
\begin{frame}
	O objetivo é percorrer exatamente uma vez todas as sete pontes da cidade que 
	conectam as duas ilhas entre si e as margens do rio, voltando ao ponto de
	partida. Existe tal rota?
	\begin{center}
		\includegraphics<1 | handout:0>[height=13\baselineskip]{Koenigsberg_Map_by_Merian-Erben_1652}
		\uncover<2->{
		\begin{tikzpicture}[every node/.style={circle,draw,blueUnB}, node distance=2cm]
			\node (A) {A};
			\node[above of=A] (B) {B} edge (A);
			\node[below of=A] (C) {C} edge (A);
			\node[right of=A, right=2cm] (D) {D} edge (A) edge (B) edge (C);
			\draw (A)..controls +(north west:1) and +(south west:1) .. (B);
			\draw (A)..controls +(south west:0.5) and +(north west:1) .. (C);
		\end{tikzpicture}}
	\end{center}
\end{frame}

\begin{frame}[label=notas]
	\begin{itemize}[<+->]
		\item O problema do inspetor de estradas foi estudado pela primeira vez 
		por Leonhard Euler (1707-1783).
		\item Euler mostrou a não existência de tal circuito através de um 
		argumento simples: 
		\subitem{\emph{Consideremos, por exemplo, a ilha da direita ($D$). Um 
		circuito qualquer deve chegar à ilha e sair dela o mesmo número de vezes. 
		Logo, para que exista um circuito euleriano, deve haver um 
		\alert<3>{número par} de pontes com extremidade nesta ilha.}}
	\end{itemize}
	\vfill
	\uncover<4->{
	\begin{block}{Teorema de Euler}
		Um grafo admite um circuito euleriano se, e somente se, o número de 
		arestas em cada vértice for par.
	\end{block}}
\end{frame}

\newcommand{\pent}[3][1]{%
	\node (#21) at(72/#1:#3cm){};
	\node (#22) at(72/#1+72:#3cm) {};
	\node (#23) at(72/#1+2*72:#3cm) {};
	\node (#24) at(72/#1+3*72:#3cm) {};
	\node (#25) at(72/#1+4*72:#3cm) {} ;
}%
\newcommand{\hpath}[1][black] {
	\draw[#1] (p11) -- (p12) -- (p13) -- (p14) -- (p15) -- (p25) -- (p35) -- (p24) -- (p34) -- (p23) -- (p33) -- (p22) -- (p32) -- (p42) -- (p43) -- (p44) -- (p45) -- (p41) -- (p31) -- (p21) -- (p11);
}

\subsection{Caixeiro viajante}
\begin{frame}[label=notas]
	Um representante de vendas de uma companhia deseja planejar uma rota, 
	iniciando na capital, na qual ele visite cada cidade exatamente uma vez, 
	voltando ao ponto de partida. Existe tal rota?
	\pause	
	\begin{center}
		\uncover<2->{
		\begin{tikzpicture}[every node/.style={circle,draw,blueUnB}, scale=.7]%, node distance=2cm]
			\pent{p1}{1}
			\pent{p2}{2}
			\pent[2]{p3}{3}
			\pent[2]{p4}{4.5}
			\hpath
			\draw<2> (p11) -- (p15);
			\draw<2> (p12) -- (p22);
			\draw<2> (p13) -- (p23);
			\draw<2> (p14) -- (p24);
			\draw<2> (p21) -- (p32);
			\draw<2> (p25) -- (p31);
			\draw<2> (p33) -- (p43);
			\draw<2> (p34) -- (p44);
			\draw<2> (p35) -- (p45);
			\draw<2> (p41) -- (p42);
			\uncover<3>{
			\hpath[red, line width=2pt]}
		\end{tikzpicture}}
	\end{center}
\end{frame}

% \begin{frame}[label=notas]
% 	\begin{itemize}[<+->]
% 		\item É o caminho Hamiltoniano (em homenagem a William Hamilton).
% 		\item É fácil convencer alguém da existência de um circuito hamiltoniano 
% 		em um grafo: basta exibir tal caminho.
% 		\item É difícil convencer alguém da não existência de um tal circuito.
% 		\item Não há método \alert<4>{eficiente} para determinar se um tal 
% 		circuito existe em um grafo.
% 		\item Existe um procedimento para verificar se um determinado grafo G 
% 		possui um circuito hamiltoniano:
% 		\subitem{Basta gerar \alert<6>{todas} as permutações circulares dos nós 
% 		deste grafo e testar para ver se uma delas corresponde a um percurso no grafo.}
% 	\end{itemize}
% \end{frame}

\begin{frame}
	\begin{itemize}[<+->]
		\item Suponha um grafo completo com 50 vértices.
		\item São 49! circuitos.
			\begin{itemize}
				\item $49! = 1{\cdot}2{\cdot}3{\cdot}4{\cdot}5{\cdot}6{\cdot}7{\cdot}8{\cdot}9{\cdots}16{\cdot}17{\cdots}32{\cdot}33{\cdots}49$
				\item $49! > 1{\cdot}2{\cdot}2{\cdot}4{\cdot}4{\cdot}4{\cdot}4{\cdot}8{\cdot}8{\cdots}16{\cdot}16{\cdots}32{\cdot}32{\cdots}32$
				\item $49! > 1{\cdot}2^2{\cdot}4^4{\cdot}8^8{\cdot}16^{16}{\cdot}32^{18}$
				\item $49! > (2^2){\cdot}(2^2)^4{\cdot}(2^3)^8{\cdot}(2^4)^{16}{\cdot}(2^5)^{18}$
				\item $49! > 2^2{\cdot}2^8{\cdot}2^{24}{\cdot}2^{64}{\cdot}2^{90}$
				\item $49! > 2^{188} =  2^8{\cdot}(2^{10})^{18} >  2^8{\cdot}(10^3)^{18}$
				\item $49! >2^8{\cdot}10^{54}$
			\end{itemize}
		\item Se um computador moderno pode realizar 200 milhões de operações 
		por segundo e a cada operação ele testar um circuito:
			\subitem[]{$$\frac{2^8{\cdot}10^{54}}{(200{\cdot}10^6){\cdot}(31.6{\cdot}10^6)}\frac{\text{circuitos}}{\frac{\text{circuitos}}{\text{segundo}}\frac{\text{segundo}}{\text{ano}}} \approx 4{\cdot}10^{40}\text{ anos.}$$}
		\item A idade do universo é da ordem de $10^{15}$ anos...
	\end{itemize}
\end{frame}